\scanpage{7}
\mspage{6}

\uline{Ex \deleted{\unsure{1}}2}\quad Soit $\deleted{\unsure{S|E}}$ $\mathcal L$ un ens ordonné,
on suppose $\mathcal L$ filtrant croissant, filtrant
décroissant, sans plus grand ni
plus petit \unsure{élémt|élément}, loc. filtrant croissant
et filtrant décroissant, divisible.

On appelera un tel ensemble une
\uline{ligne ordonnée}

\[
S_0\subset P_2(\mathcal L)\qquad
\deleted{\text{Soit }\{x,y\}\in P_2(\mathcal L)}
\]

\textquotedbl\quad \unsure{partie|parties} \deleted{qui} de card. 2 \deleted{\unsure{une}} qui sont \uline{tot.} ordonnées

\[
\begin{aligned}
&\mathrm{Modf}(\mathcal L)=\text{parties finies tot ordonnées}\\
&\hspace{8em}\deleted{\text{i.e. telles que}}
\end{aligned}
\]

\[
\begin{aligned}
 \text{Soit }T\in\mathrm{Modf}(\mathcal L),\qquad &T=\{t_{\unsure{p|1}},t_2,\ldots,t_{n}\}\\
 &\quad t_1<\unsure{t|l|1}_2<\ldots<t_n
\end{aligned}
\]

Alors les composantes connexes sont

\[
\underbrace{\mathopen{]}-\infty,t_1\mathclose{[}}_{t_1},\
\underbrace{\mathopen{]}t_1,t_2\mathclose{[}}_{t_1,t_2},\
\ldots\quad
\underbrace{\mathopen{]}t_{n-1},t_n\mathclose{[}}_{t_{n-1}\,t_n},\
\underbrace{\mathopen{]}t_n,\infty\mathclose{]}}_{t_n}
\]

% slanted
\margin{Pts frontière}

($n+1$ composantes connexes)

% paper watermark; mirror impression across the letters
\rotatebox{90}{\textsc{\unsure{t}rong Extra Str}}
